% Source: https://artofproblemsolving.com/community/c3h3579385_1434th_post
% Stable identity: aops-topic-3579385-q14
Write the original tangent lengths from $A,B,C$ as $x,y,z$, respectively, so $AB=x+y$, $AC=x+z$, $BC=y+z$, and $BE=z$, $EC=y$. If $AE=t$, the two smaller semiperimeters both equal $\dfrac{x+y+z+t}{2}$. Their area ratio is therefore $4:11$, and since they have a common altitude to $BC$, $z:y=4:11$. Let the original inradius be $r$. Place $B=(0,0)$, $C=(y+z,0)$, $E=(z,0)$ with $A$ above the axis. Let the two smaller incentres be $I,J$. Since $\cot\dfrac{B}{2}=\dfrac{y}{r}$ and $\cot\dfrac{C}{2}=\dfrac{z}{r}$, their coordinates are $I=(\dfrac{4y}{r},4)$ and $J=(y+z-\dfrac{11z}{r},11)$. The incircle of $ACE$ lies inside $ABC$, so the original triangle's maximum inscribed radius gives $r\ge11$. Equality would require the unique incircle of $ABC$ to lie inside $ACE$, but its tangency point on $AB$ lies outside $ACE$. Thus $r>11$. Relative to $E$, write $I=(-u,4)$ and $J=(v,11)$, with $u=z-\dfrac{4y}{r}>0$ and $v=y-\dfrac{11z}{r}>0$. The angles at $E$ of the two subtriangles sum to $180^\circ$, so their internal bisectors $EI,EJ$ are perpendicular. Hence $uv=44$. Also $IJ=25$ and the vertical difference is $7$, so $u+v=\sqrt{625-49}=24$. Thus $(u,v)$ is either $(2,22)$ or $(22,2)$. Using $z=\dfrac{4y}{11}$ gives\[\dfrac{u}{v}=\dfrac{\dfrac{4}{11}-\dfrac{4}{r}}{1-\dfrac{4}{r}}.\]For $(u,v)=(2,22)$ this yields $r=\dfrac{40}{3}$; for $(22,2)$ it yields $r=\dfrac{440}{117}<11$, impossible. Therefore $\boxed{r=\dfrac{40}{3}}$.
\begin{asy}
import geometry;
size(260);
pair B=(0,0),C=(300.0/7,0),A=(1786.0/35,264.0/5),D=(220.0/7,0),E=(80.0/7,0),O=(220.0/7,40.0/3),I=(66.0/7,4),J=(234.0/7,11);draw(A--B--C--cycle);draw(A--E);draw(circle(O,40.0/3),gray);draw(circle(I,4),blue);draw(circle(J,11),red);draw(E--I--J--E,dashed);label("$A$",A,N);label("$B$",B,SW);label("$C$",C,SE);label("$D$",D,S);label("$E$",E,S);label("$O$",O,N);label("$I$",I,W);label("$J$",J,E);
\end{asy}
